\chapter{Couche 3 — La session et le navigateur}

L'utilisateur est authentifié. Reste à empêcher qu'on lui vole sa session, ou
qu'on la lui fasse utiliser à son insu.

% -----------------------------------------------------------------------------
\mesure{Le cookie de session}{mesure 5 · \fichier{config.py}}

\menace
Après la connexion, c'est le cookie qui prouve l'identité. Le voler, c'est
devenir l'utilisateur sans jamais connaître son mot de passe.

\parade
Quatre attributs, chacun fermant une porte différente.

\danslecode
\begin{lstlisting}[language=Python]
SESSION_COOKIE_HTTPONLY   = True             # inaccessible a JavaScript
SESSION_COOKIE_SAMESITE   = "Strict"         # jamais envoye depuis un autre site
SESSION_COOKIE_SECURE     = HTTPS            # transmis seulement en TLS
SESSION_COOKIE_NAME       = "insacloud_session"
PERMANENT_SESSION_LIFETIME = timedelta(hours=2)
\end{lstlisting}

\begin{center}
\small
\begin{tabularx}{\textwidth}{@{}p{3.4cm}X@{}}
\toprule
\textbf{Attribut} & \textbf{La porte qu'il ferme} \\
\midrule
\cmd{HttpOnly} & Un script injecté dans la page ne peut pas lire le cookie :
une faille XSS ne se transforme pas en vol de session. \\
\cmd{SameSite=Strict} & Le navigateur n'envoie pas le cookie quand la requête
vient d'un autre site. C'est la défense structurelle contre le CSRF, celle qui
agit avant même que la requête parte. \\
\cmd{Secure} & Le cookie ne circule jamais en clair. Activé dès que le site est
en HTTPS --- désactivé en développement local, sinon le cookie ne serait jamais
émis et rien ne fonctionnerait. \\
\cmd{2 heures} & Une session oubliée sur un poste de salle TP finit par
expirer. \\
\bottomrule
\end{tabularx}
\end{center}

\begin{piege}{Une session Flask est signée, pas chiffrée}
C'est la méprise la plus répandue. Le contenu du cookie est encodé en base64 :
n'importe qui peut le \emph{lire}. Ce que la clé secrète garantit, c'est qu'on
ne peut pas le \emph{modifier} sans invalider la signature.

Conséquence directe sur le code : la session ne contient que
l'\textbf{identifiant} de l'utilisateur. Jamais son rôle, jamais un jeton,
jamais un mot de passe. Tout le reste est relu en base à chaque requête.
\end{piege}

\limite
Un cookie volé reste valable jusqu'à son expiration : il n'est lié ni à
l'adresse IP ni au navigateur. Lier une session à son adresse casserait les
connexions mobiles et n'arrêterait qu'un attaquant peu motivé.

% -----------------------------------------------------------------------------
\mesure{Les jetons CSRF}{mesure 6 · tous les formulaires}

\menace
Un site tiers contient une image ou un formulaire qui pointe vers
\cmd{/instances/42/supprimer}. L'utilisateur, connecté à InsaCloud dans un
autre onglet, le visite : son navigateur joint son cookie, et la requête
part avec son identité.

\parade
Chaque formulaire porte un jeton aléatoire lié à la session ; le serveur refuse
toute requête qui modifie l'état sans ce jeton. Un site tiers ne peut pas le
deviner --- il ne peut pas lire la page qui le contient, la politique de même
origine l'en empêche.

\danslecode
\cmd{WTF\_CSRF\_ENABLED = True}, y compris en configuration de test : la
protection CSRF reste donc exercée par la suite de tests, et une régression qui
la désactiverait ferait échouer la chaîne d'intégration.

\begin{retenir}
\cmd{SameSite=Strict} et le jeton CSRF font le même travail à deux endroits
différents --- le navigateur et le serveur. On garde les deux : le premier
dépend du navigateur de l'utilisateur, le second de personne.
\end{retenir}

\limite
Le CSRF ne protège pas d'un script qui s'exécute \emph{dans} la page : celui-là
lit le jeton. C'est la politique de contenu qui l'en empêche.

% -----------------------------------------------------------------------------
\mesure{Les en-têtes de sécurité}{mesure 7 · \fichier{app.py}}

\menace
Le navigateur est permissif par défaut : il exécute les scripts qu'on lui
donne, affiche la page dans un cadre invisible, devine les types de contenu,
et transmet l'URL consultée au site suivant. Chacun de ces comportements est
une attaque possible.

\parade
Huit en-têtes posés sur \emph{chaque} réponse, par un \cmd{after\_request}
unique --- impossible d'en oublier un sur une route.

\danslecode
\begin{lstlisting}[language=Python]
@app.after_request
def security_headers(response):
    csp = ("default-src 'self'; script-src 'self'; style-src 'self'; "
           "img-src 'self' data:; font-src 'self'; connect-src 'self'; "
           "object-src 'none'; base-uri 'self'; form-action 'self'; "
           "frame-ancestors 'none'")
    if HTTPS:
        csp += "; upgrade-insecure-requests"
        response.headers["Strict-Transport-Security"] = \
            "max-age=31536000; includeSubDomains"
    response.headers["Content-Security-Policy"]   = csp
    response.headers["X-Content-Type-Options"]    = "nosniff"
    response.headers["X-Frame-Options"]           = "DENY"
    response.headers["Referrer-Policy"]           = "no-referrer"
    response.headers["Permissions-Policy"]        = \
        "camera=(), microphone=(), geolocation=(), payment=()"
    response.headers["Cross-Origin-Opener-Policy"]   = "same-origin"
    response.headers["Cross-Origin-Resource-Policy"] = "same-origin"
\end{lstlisting}

\begin{center}
\small
\begin{tabularx}{\textwidth}{@{}p{4.3cm}X@{}}
\toprule
\textbf{En-tête} & \textbf{Ce qu'il interdit} \\
\midrule
\cmd{Content-Security-Policy} & Charger un script, un style, une police ou une
image d'ailleurs que du site lui-même. \cmd{object-src 'none'} bloque les
greffons, \cmd{form-action 'self'} empêche un formulaire de poster ailleurs,
\cmd{base-uri 'self'} empêche de réécrire la base des liens relatifs. \\
\cmd{frame-ancestors 'none'} \newline \cmd{X-Frame-Options: DENY} & Afficher le
site dans un cadre sur un autre site --- la parade au \emph{clickjacking}. Les
deux en-têtes disent la même chose ; le second pour les navigateurs anciens. \\
\cmd{X-Content-Type-Options} & Deviner le type d'un fichier plutôt que de
croire l'en-tête : un fichier téléversé ne peut pas être réinterprété en
script. \\
\cmd{Referrer-Policy: no-referrer} & Transmettre au site suivant l'URL que
l'utilisateur consultait --- qui contient parfois un identifiant. \\
\cmd{Permissions-Policy} & Demander la caméra, le micro, la position ou le
paiement. La plateforme n'en a aucun usage ; autant les refuser explicitement. \\
\cmd{Strict-Transport-Security} & Revenir en HTTP pendant un an, y compris sur
les sous-domaines. Posé seulement si le site est réellement en TLS --- le poser
en HTTP local rendrait le site inaccessible. \\
\cmd{Cross-Origin-*} & Partager le contexte de navigation ou les ressources
avec une autre origine. \\
\bottomrule
\end{tabularx}
\end{center}

\begin{piege}{\cmd{default-src 'self'} interdit les CDN --- et c'est voulu}
Une politique stricte a un coût concret : on ne peut plus écrire
\cmd{<script src="https://cdn.tailwindcss.com">}. Tailwind est donc
\textbf{précompilé} et servi depuis le site.

Le gain dépasse la sécurité : la plateforme ne dépend plus d'un service
extérieur qui pourrait tomber, changer de contenu ou observer les visiteurs.
Une politique de contenu qu'on affaiblit pour faire marcher un CDN ne protège
plus de grand-chose --- autant la tenir.
\end{piege}

\limite
La CSP limite les dégâts d'une injection, elle ne l'empêche pas. La vraie
défense contre le XSS reste l'échappement automatique des gabarits Jinja2, que
la CSP double en dernier recours.
